\chapter{\MakeUppercase{Technical Specification}}


% -----------------------------
% Table 3.1
% -----------------------------
\section{Requirements}
\subsection{Functional and Non-Functional Requirements (Level 2)}
\paragraph{}
The system requirements define the core capabilities and operational constraints of the Secure Messaging App. 
\begin{itemize}
	\item \textbf{Functional Requirements:} 
	\begin{itemize}
		\item \textit{User Authentication:} Secure identity verification and session management via JWT.
		\item \textit{Key Generation \& Distribution:} Generation of classical (ECC) and post-quantum (ML-KEM) key pairs natively on the client.
		\item \textit{Secure Messaging:} Real-time text transmission strictly over WebSockets using AES-GCM-256 authenticated encryption.
		\item \textit{Local Storage:} Secure handling and querying of user credentials and historical messages entirely offline via IndexedDB.
	\end{itemize}
	\item \textbf{Non-Functional Requirements:}
	\begin{itemize}
		\item \textit{Security Requirements:} Perfect forward secrecy, post-quantum resilience, and a zero-knowledge trust model extending to the relay infrastructure.
		\item \textit{Performance Constraints:} Key derivation and message encryption speeds must stay within acceptable limits for a real-time chat feel despite using heavier lattice-based cryptography.
	\end{itemize}
\end{itemize}

\section{Feasibility Study}
\paragraph{Usage of Symbols} The symbols: \ac{alpha}, \ac{beta}, \ac{delta}, etc., can be defined with the acronym section and used in the text as \verb|\ac{alpha}, \ac{beta}, \ac{delta}|, respectively. 

\paragraph{}
A comprehensive feasibility study was undertaken to evaluate the integration of standard, robust Post-Quantum Cryptography into an Electron-based web stack:
\begin{itemize}
	\item \textbf{Technical Feasibility:} Utilizing a pure-TypeScript implementation of FIPS 203 (\textit{mlkem}) alongside the browser's native WebCrypto API ensures that highly intensive operations (like Kyber encapsulation/decapsulation and NIST P-384 ECDH) are fully supported without native bindings, minimizing cross-platform build friction.
	\item \textbf{Operational Feasibility:} By wrapping the React interface within Electron, standard desktop functionality (notifications, sandboxed execution, IPC) operates effectively while masking complex cryptographic states from the user.
\end{itemize}


\section{System Specification}
\paragraph{}
The overarching architecture consists of highly independent client and server specifications operating in a decentralized trust arrangement. The relay server is conceptually "blind", routing binary blobs containing ciphertexts.

\begin{itemize}
	\item \textbf{Client Node:} Encapsulated React application maintaining an isolated IndexedDB storage for state. All key agreement derivations (Hybrid X3DH via HKDF-SHA-384) are processed within a hardened Electron renderer process leveraging strict CSPs.
	\item \textbf{Relay Server:} A high-throughput Fastify (Node.js) application exposing both an initial REST API (for polling prekey bundles) and an active WebSocket protocol (for immediate message routing). SQLite handles temporary persistent state like undelivered messages or offline clients. 
\end{itemize}

\section{Preparing your report/thesis/article using \LaTeX}
\paragraph{}
\subsection{Tables}
In the report, table can be defined in multiple ways according to the need.  For example, the latex code
{\scriptsize
\begin{verbatim}
\begin{table}[h!]
	\centering
	\caption{Country list}
	\renewcommand{\arraystretch}{1.3} % Row height
	\begin{tabular}{|l|l|l|l|}
		\hline
		\textbf{Country/Area Name} & \textbf{Description} & \textbf{ISO alpha 3 Code} & \textbf{ISO numeric Code} \\
		\hline
		Afghanistan & It is a Country & AFG & 004 \\
		Aland Islands & It is an Island & ALA & 248 \\
		Albania & It is a small city & ALB & 008 \\
		Algeria & It is a Country & DZA & 012 \\
		American Samoa & It is a Country & ASM & 016 \\
		\hline
	\end{tabular}
\end{table}
\end{verbatim}
}
defines the simple Table 3.1 in the document which is restricted to a single page.
\begin{table}[h!]
\centering
\caption{Country list}
\renewcommand{\arraystretch}{1.3} % Row height
\begin{tabular}{|l|l|l|l|}
\hline
\textbf{Country/Area Name} & \textbf{Description} & \textbf{ISO alpha 3 Code} & \textbf{ISO numeric Code} \\
\hline
Afghanistan & It is a Country & AFG & 004 \\
Aland Islands & It is an Island & ALA & 248 \\
Albania & It is a small city & ALB & 008 \\
Algeria & It is a Country & DZA & 012 \\
American Samoa & It is a Country & ASM & 016 \\
\hline
\end{tabular}
\end{table}

\pagebreak

The oversized table can be fit with in the page as follows: 
{\scriptsize
	\begin{verbatim}
		\begin{table}[h!]
			\centering
			\caption{Data units, sources, and dates}
			\renewcommand{\arraystretch}{1.3}
			\resizebox{\textwidth}{!}{
				\begin{tabular}{|l|l|l|l|}
					\hline
					\textbf{Variable} & \textbf{Dates} & \textbf{Units} & \textbf{Source} \\
					\hline
					Nominal Physical Capital Stock & 1950-1990 & Billions US\$ & Nehru and Dhareshwar (1993) \\
					Total Population & 1950-1990 & Billions & Nehru and Dhareshwar (1993) \\
					Nominal GDP & 1950-1990 & Billions US\$ & PWT \\
					Real GDP per capita & 1950-1990 & 2005 US\$ per capita & PWT \\
					\hline
			\end{tabular}}
		\end{table}
	\end{verbatim}
}
which produces the Table 3.2.

% -----------------------------
% Table 3.2
% -----------------------------
\begin{table}[h]
\centering
\caption{Data units, sources, and dates}
\renewcommand{\arraystretch}{1.3}
\resizebox{\textwidth}{!}{
\begin{tabular}{|l|l|l|l|}
\hline
\textbf{Variable} & \textbf{Dates} & \textbf{Units} & \textbf{Source} \\
\hline
Nominal Physical Capital Stock & 1950-1990 & Billions US\$ & Nehru and Dhareshwar (1993) \\
Total Population & 1950-1990 & Billions & Nehru and Dhareshwar (1993) \\
Nominal GDP & 1950-1990 & Billions US\$ & PWT \\
Real GDP per capita & 1950-1990 & 2005 US\$ per capita & PWT \\
\hline
\end{tabular}}
\end{table}

\paragraph{Defining the table span for more than one page}. 
The longtable (span for more than one page) can be achieved using following code. 
{\scriptsize
	\begin{verbatim}
		\begin{longtable}{|l|l|l|}
			\caption{This table can be used if more than one page table content is required. However it is limited to maximum of two pages only.} \label{tab:long} \\
			
			\hline \multicolumn{1}{|c|}{\textbf{First column}} & \multicolumn{1}{c|}{\textbf{Second column}} & \multicolumn{1}{c|}{\textbf{Third column}} \\ \hline 
			\endfirsthead
			
			\multicolumn{3}{c}%
			{{\bfseries \tablename\ \thetable{} -- continued from previous page}} \\
			\hline \multicolumn{1}{|c|}{\textbf{First column}} & \multicolumn{1}{c|}{\textbf{Second column}} & \multicolumn{1}{c|}{\textbf{Third column}} \\ \hline 
			\endhead
			
			\hline \multicolumn{3}{|r|}{{Continued on next page}} \\ \hline
			\endfoot
			
			\hline \hline
			\endlastfoot
			One & abcdef ghjijklmn & 123.456778 \\
			One & abcdef ghjijklmn & 123.456778 \\
			....
			....
			....
			One & abcdef ghjijklmn & 123.456778 \\
			One & abcdef ghjijklmn & 123.456778 \\
		\end{longtable}
	\end{verbatim}
}

\begin{longtable}{|l|l|l|}
	\caption{This table can be used if more than one page table content is required. However it is limited to maximum of two pages only.} \label{tab:long} \\
	
	\hline \multicolumn{1}{|c|}{\textbf{First column}} & \multicolumn{1}{c|}{\textbf{Second column}} & \multicolumn{1}{c|}{\textbf{Third column}} \\ \hline 
	\endfirsthead
	
	\multicolumn{3}{c}%
	{{\bfseries \tablename\ \thetable{} -- continued from previous page}} \\
	\hline \multicolumn{1}{|c|}{\textbf{First column}} & \multicolumn{1}{c|}{\textbf{Second column}} & \multicolumn{1}{c|}{\textbf{Third column}} \\ \hline 
	\endhead
	
	\hline \multicolumn{3}{|r|}{{Continued on next page}} \\ \hline
	\endfoot
	
	\hline \hline
	\endlastfoot
	
	One & abcdef ghjijklmn & 123.456778 \\
	One & abcdef ghjijklmn & 123.456778 \\
	One & abcdef ghjijklmn & 123.456778 \\
	One & abcdef ghjijklmn & 123.456778 \\
	One & abcdef ghjijklmn & 123.456778 \\
	One & abcdef ghjijklmn & 123.456778 \\
	One & abcdef ghjijklmn & 123.456778 \\
	One & abcdef ghjijklmn & 123.456778 \\
	One & abcdef ghjijklmn & 123.456778 \\
	One & abcdef ghjijklmn & 123.456778 \\
	One & abcdef ghjijklmn & 123.456778 \\
	One & abcdef ghjijklmn & 123.456778 \\
	One & abcdef ghjijklmn & 123.456778 \\
	One & abcdef ghjijklmn & 123.456778 \\
	One & abcdef ghjijklmn & 123.456778 \\
	One & abcdef ghjijklmn & 123.456778 \\
	One & abcdef ghjijklmn & 123.456778 \\
	One & abcdef ghjijklmn & 123.456778 \\
	One & abcdef ghjijklmn & 123.456778 \\
	One & abcdef ghjijklmn & 123.456778 \\
	One & abcdef ghjijklmn & 123.456778 \\
	One & abcdef ghjijklmn & 123.456778 \\
	One & abcdef ghjijklmn & 123.456778 \\
	One & abcdef ghjijklmn & 123.456778 \\
	One & abcdef ghjijklmn & 123.456778 \\
	One & abcdef ghjijklmn & 123.456778 \\
	One & abcdef ghjijklmn & 123.456778 \\
	One & abcdef ghjijklmn & 123.456778 \\
	One & abcdef ghjijklmn & 123.456778 \\
	One & abcdef ghjijklmn & 123.456778 \\
	One & abcdef ghjijklmn & 123.456778 \\
	One & abcdef ghjijklmn & 123.456778 \\
	One & abcdef ghjijklmn & 123.456778 \\
	One & abcdef ghjijklmn & 123.456778 \\
	One & abcdef ghjijklmn & 123.456778 \\
	One & abcdef ghjijklmn & 123.456778 \\
	One & abcdef ghjijklmn & 123.456778 \\
	One & abcdef ghjijklmn & 123.456778 \\
	One & abcdef ghjijklmn & 123.456778 \\
	One & abcdef ghjijklmn & 123.456778 \\
	One & abcdef ghjijklmn & 123.456778 \\
	One & abcdef ghjijklmn & 123.456778 \\
	One & abcdef ghjijklmn & 123.456778 \\
	One & abcdef ghjijklmn & 123.456778 \\
	One & abcdef ghjijklmn & 123.456778 \\
	One & abcdef ghjijklmn & 123.456778 \\
	One & abcdef ghjijklmn & 123.456778 \\
	One & abcdef ghjijklmn & 123.456778 \\
	One & abcdef ghjijklmn & 123.456778 \\
	One & abcdef ghjijklmn & 123.456778 \\
	One & abcdef ghjijklmn & 123.456778 \\
	One & abcdef ghjijklmn & 123.456778 \\
	One & abcdef ghjijklmn & 123.456778 \\
	One & abcdef ghjijklmn & 123.456778 \\
	One & abcdef ghjijklmn & 123.456778 \\
	One & abcdef ghjijklmn & 123.456778 \\
	One & abcdef ghjijklmn & 123.456778 \\
	One & abcdef ghjijklmn & 123.456778 \\
	One & abcdef ghjijklmn & 123.456778 \\
	One & abcdef ghjijklmn & 123.456778 \\
	One & abcdef ghjijklmn & 123.456778 \\
	One & abcdef ghjijklmn & 123.456778 \\
	One & abcdef ghjijklmn & 123.456778 \\
	One & abcdef ghjijklmn & 123.456778 \\
	One & abcdef ghjijklmn & 123.456778 \\
	One & abcdef ghjijklmn & 123.456778 \\
	One & abcdef ghjijklmn & 123.456778 \\
	One & abcdef ghjijklmn & 123.456778 \\
	One & abcdef ghjijklmn & 123.456778 \\
	One & abcdef ghjijklmn & 123.456778 \\
	One & abcdef ghjijklmn & 123.456778 \\
	One & abcdef ghjijklmn & 123.456778 \\
	One & abcdef ghjijklmn & 123.456778 \\
	One & abcdef ghjijklmn & 123.456778 \\
	One & abcdef ghjijklmn & 123.456778 \\
	One & abcdef ghjijklmn & 123.456778 \\
	One & abcdef ghjijklmn & 123.456778 \\
	One & abcdef ghjijklmn & 123.456778 \\
	One & abcdef ghjijklmn & 123.456778 \\
	One & abcdef ghjijklmn & 123.456778 \\
	One & abcdef ghjijklmn & 123.456778 \\
	One & abcdef ghjijklmn & 123.456778 \\
	One & abcdef ghjijklmn & 123.456778 \\
\end{longtable}

The table can also be presented in the Landscape orientation as follows: 
\begin{sidewaystable}
	\centering
	\caption{This is long table presented in Landscape orientation}
	\begin{tabular}{|c|c|c|c|c|c|c|c|c|c|c|c}
		\hline
		col 1 & col 2 & col 3 & col 4 & col 5 & col 6 & col 7 & col 8 & col 9 & col 10 & col 11 & col 12 \\
		One & abcdef ghjijklmn & 123.456778 & One & abcdef ghjijklmn & 123.456778 & One & abcdef ghjijklmn & 123.456778 & One & abcdef ghjijklmn &  123.456778 \\
		One & abcdef ghjijklmn & 123.456778 & One & abcdef ghjijklmn & 123.456778 & One & abcdef ghjijklmn & 123.456778 & One & abcdef ghjijklmn &  123.456778 \\
		One & abcdef ghjijklmn & 123.456778 & One & abcdef ghjijklmn & 123.456778 & One & abcdef ghjijklmn & 123.456778 & One & abcdef ghjijklmn &  123.456778 \\
		One & abcdef ghjijklmn & 123.456778 & One & abcdef ghjijklmn & 123.456778 & One & abcdef ghjijklmn & 123.456778 & One & abcdef ghjijklmn &  123.456778 \\
		\hline
	\end{tabular}
\end{sidewaystable}	


\subsection{Equations}
In line equation can be defined in latex as \verb|$\hat{y} = f(x; W, b)$|. This code produce the formula as $\hat{y} = f(x; W, b)$ . 

\paragraph{Equations}
In latex, the equation can be defined using following commands: 
\begin{verbatim}
	\begin{equation}
		\delta^{L} = \frac{\partial \mathcal{L}}{\partial a^{L}} \odot \sigma'(z^{L})
	\end{equation}
\end{verbatim}
which will produce the Equation 3.1. You can use \verb|$\delta^{L}$| to represent $\delta^{L}$ in the running text.  For example, $\odot$ is a Hadamard operator. 
\begin{equation}
	\delta^{L} = \frac{\partial \mathcal{L}}{\partial a^{L}} \odot \sigma'(z^{L})
\end{equation}

\subsection{Figures}
\paragraph{Figure} Figures can be included in the latex document. One simple figure/image can be included as 
\begin{verbatim}
	\begin{figure}[h]
		\centering
		\includegraphics[width=0.5\linewidth]{images/sample-graph.pdf}
		\caption{Sample Diagram}
		\label{fig:universe}
	\end{figure}
\end{verbatim} 
which will produce the output as 
\begin{figure}[!h]
	\centering
	\includegraphics[width=0.5\linewidth]{images/sample-graph.pdf}
	\caption{Sample Diagram}
	\label{fig:universe1}
\end{figure}
\linebreak
More than one figures can be shown adjacent using following code
\begin{verbatim}
	\begin{figure}[h!]
		\centering
		% Left image (a)
		\begin{subfigure}[b]{0.45\linewidth}
			\centering
			\includegraphics[width=\textwidth]{images/sample-graph.pdf}  % replace with your image filename
			\caption{Example Graph}
			\label{fig:fruitsA}
		\end{subfigure}
		\hfill
		% Right image (b)
		\begin{subfigure}[b]{0.45\linewidth}
			\centering
			\includegraphics[width=\textwidth]{images/sample-graph.pdf} % replace with your image filename
			\caption{Example Graph}
			\label{fig:fruitsC}
		\end{subfigure}
		
		\caption{Commonly available fruits}
		\label{fig:common_fruits}
	\end{figure}
\end{verbatim}
The above code will produce the results as 
\begin{figure}[h!]
	\centering
	% Left image (a)
	\begin{subfigure}[b]{0.45\linewidth}
		\centering
		\rotatebox{45}{
		\includegraphics[width=\textwidth]{images/sample-graph.pdf} %	 replace with your image filename
		}
		
	
		\label{fig:fruitsAA}
	\end{subfigure}
	\hfill
	% Right image (b)
	\begin{subfigure}[b]{0.45\linewidth}
		\centering
		\rotatebox{-45}{
		\includegraphics[width=\textwidth]{images/sample-graph.pdf} % replace with your image filename
		}
		\label{fig:fruitsCC}
	\end{subfigure}
	
	\caption{Figures with rotation}
	\label{fig:common_fruitsA}
\end{figure}

\subsection{Pseudo code / Algorithm}
The algorithm or pseudo code can be added using following code
\begin{verbatim}
	\begin{algorithm}
		\caption{An example algorithm with caption}\label{alg:cap}
		\begin{algorithmic}
			\Require $n \geq 0$
			\Ensure $y = x^n$
			\State $y \gets 1$
			\State $X \gets x$
			\State $N \gets n$
			\While{$N \neq 0$}
			\If{$N$ is even}
			\State $X \gets X \times X$
			\State $N \gets \frac{N}{2}$  \Comment{This is a comment}
			\ElsIf{$N$ is odd}
			\State $y \gets y \times X$
			\State $N \gets N - 1$
			\EndIf
			\EndWhile
		\end{algorithmic}
	\end{algorithm}
\end{verbatim}
This code will produce the following output in the report.
\begin{algorithm}
	\caption{An example algorithm with caption}\label{alg:cap}
	\begin{algorithmic}
		\Require $n \geq 0$
		\Ensure $y = x^n$
		\State $y \gets 1$
		\State $X \gets x$
		\State $N \gets n$
		\While{$N \neq 0$}
		\If{$N$ is even}
		\State $X \gets X \times X$
		\State $N \gets \frac{N}{2}$  \Comment{This is a comment}
		\ElsIf{$N$ is odd}
		\State $y \gets y \times X$
		\State $N \gets N - 1$
		\EndIf
		\EndWhile
	\end{algorithmic}
\end{algorithm}
